% Section: Conclusion
\section{Conclusion}
This paper presented two architectural extensions to the DeepMTL2R framework, which are Dynamic Feature Gating (DFG) for suppressing low-information features, and Matryoshka Feature Projection (MFP) for generating nested, sliceable representations. However, these extensions failed to demonstrate their intended benefits due to a numerical collapse in the listwise LambdaLoss objective. Specifically, directly using large unnormalized auxiliary targets (e.g., click counts up to $13.15 \times 10^{6}$) caused the exponential gain function to exceed \texttt{float32} limits. This resulted in \texttt{NaN} gradients that froze all model weights at initialization, yielding entirely stagnant evaluation metrics, including a constant NDCG@10 $\approx 0.212$ and identical Matryoshka dimension scores.

Despite the training failure, the underlying cause has been clearly identified and can be addressed systematically. We recommend normalizing auxiliary targets on a per-slate basis (via min-max scaling or ordinal binning) and employing task-specific loss functions that separate regression-scale targets from the ranking objective. Once numerical stability is restored through these measures, future work should reevaluate DFG's feature interpretability and MFP's dimensional efficiency across the complete MSLR-WEB10K dataset.